\section{Aggregate implications}\label{sec:macro}

\subsection{Stepping on toes, measured}

Growth models since \citet{jones1995} let the flow of new ideas scale with
researchers as $N^{\lambda}$, $\lambda<1$, and call $1-\lambda$ ``stepping on
toes''. The parameter is reduced form. Here the flow of credited novel output
is $Y=N\,a\,F(M^*(N))$, because in equilibrium every researcher's career payoff
equals the frontier's. The model therefore gives $\lambda$ as an equilibrium object.

\begin{proposition}[Stepping on toes]\label{prop:sot}
In the large-topic limit,
\begin{equation}\label{eq:sot}
\lambda\equiv\frac{\dd\ln Y}{\dd\ln N}=1-\eta\,\frac{\dd\ln M^*}{\dd\ln N}=1-\mathcal E\,|\eps_w(P^*)|.
\end{equation}
If, in addition, topics opened later are less fertile, with fertility elasticity $\eta_f\ge0$ (\cref{app:fertility}), the right-hand side of \eqref{eq:sot} is an upper bound on $\lambda$.
\end{proposition}

The formula is a sufficient statistic built from two observable objects. The
first is the volume elasticity of density, $\mathcal E$ (Test 1). The second is
the elasticity of the entrant's payoff curve at equilibrium density, $\eps_w=\eps_v-\kappa P^*$. That is the
incumbent's returns gradient minus the catch-up premium (Test 4 and
\cref{prop:mature}). Crowding does not \emph{amplify} frontier depletion, as
the draft conjectured. It \emph{buffers} it. The density margin absorbs part of
the growth in researchers, so fewer new topics are opened and the frontier
deteriorates more slowly. The cost of crowding shows up in the \emph{level} of
productivity, as the dissipated rent $\hat w-F$, and not in its elasticity.
With $\mathcal E\approx0.2$ and the calibrated slope of $w$, crowding and access
depletion together imply $\lambda$ close to one (\cref{sec:quant}). The bulk of the measured
decline in research productivity \citep{bloometal2020} must come from
elsewhere: falling fertility of new topics ($\eta_f$, not identified by density
data) or selection on researcher ability, as in \citet{ekerdtwu2024}.

\subsection{A Solow paradox for science}

The last half-century brought enormous improvements in access to the
literature: online journals, preprint servers, full-text search, and now
language models that read for us. Measured research productivity fell
throughout \citep{bloometal2020}. \Cref{prop:ai}(iii) suggests one reason.

\begin{proposition}[Reading improvements are dissipated]\label{prop:solow}
Let $\omega$ be any parameter that raises $w$ at every density and leaves $F$
unchanged. In the large-topic limit, $\dd P^*/\dd\omega>0$ and
$\dd M^*/\dd\omega<0$, and
\[
\frac{\dd\ln(Y/N)}{\dd\omega}=\eta\left|\frac{\dd\ln M^*}{\dd\omega}\right|\ge0,
\]
with equality if and only if $\eta=0$. By contrast, $\dd\ln(Y/N)/\dd\ln a=1-\mathcal E|\eps_w|$.
\end{proposition}

Better access to knowledge raises the payoff of working in existing topics.
Newcomers respond by crowding them until they are no better than the frontier.
What remains is denser topics, fewer new ones, and the same output per
researcher. The prediction matches \citet{evans2008}, who finds that online
access narrowed the set of papers that get cited. It also matches
\citet{haoetal2026}, who find that AI tools ``contract science's focus''. Gains
from access to frontier knowledge measured at the level of individuals or
topics \citep{iariaetal2018,biasimoser2021} are real, but in equilibrium they
are partly competed away at the level of science.
